\documentclass[11pt,a4paper]{article}

% --- Pacchetti ---
\usepackage[utf8]{inputenc}
\usepackage[T1]{fontenc}
\usepackage[italian]{babel}
\usepackage{amsmath, amssymb}
\usepackage{booktabs}
\usepackage{geometry}
\usepackage{hyperref}
\usepackage{enumitem}
\usepackage{array}
\usepackage{multirow}
\usepackage{fancyhdr}
\usepackage{xcolor}
\usepackage{tabularx}

\geometry{margin=2.5cm}
\hypersetup{colorlinks=true, linkcolor=blue!60!black, urlcolor=blue!60!black}

\pagestyle{fancy}
\fancyhf{}
\rhead{MARL Solar Edge -- Documento di Progetto}
\lhead{Tesi Magistrale}
\cfoot{\thepage}

\title{\textbf{Sistema Multi-Agente per il Bilanciamento Energetico\\di un Cluster Edge con Energy Harvesting Fotovoltaico}\\[0.5em]\large Architettura, Modello e Metriche di Valutazione}
\author{Alessandro Colantuoni}
\date{\today}

\begin{document}
\maketitle
\tableofcontents
\newpage

% ======================================================================
\section{Introduzione e Motivazione}
% ======================================================================

Il \emph{Multi-access Edge Computing} (MEC) avvicina le risorse computazionali ai dispositivi IoT, riducendo latenza e consumo di banda. In scenari outdoor (smart agriculture, monitoraggio ambientale, smart city), i nodi edge sono alimentati a batteria e dotati di pannelli fotovoltaici per l'\emph{Energy Harvesting}.

Il problema centrale di questa tesi \`e il \textbf{Task Offloading} con vincoli energetici: dato un flusso continuo di task computazionali eterogenei, come distribuirli su un cluster di nodi edge in modo da:
\begin{enumerate}[nosep]
    \item Massimizzare il throughput (task completati localmente);
    \item Mantenere bilanciata l'energia dell'intero cluster;
    \item Sfruttare la \emph{Solar Wave}: il fatto che pannelli con orientamenti diversi ricevono picchi solari in momenti diversi della giornata.
\end{enumerate}

L'approccio adottato \`e il \textbf{Multi-Agent Reinforcement Learning (MARL)} con l'algoritmo \textbf{IPPO} (Independent Proximal Policy Optimization): ogni nodo possiede una policy neurale indipendente che impara autonomamente, senza comunicazione esplicita con gli altri nodi.

% ======================================================================
\section{Architettura del Sistema}
% ======================================================================

\subsection{Componenti Software}

Il sistema \`e costruito su tre framework:
\begin{itemize}[nosep]
    \item \textbf{PettingZoo} (v1.24+): definisce l'ambiente multi-agente con API \texttt{ParallelEnv}, dove tutti gli agenti agiscono simultaneamente ad ogni step.
    \item \textbf{Ray RLlib} (v2.x): orchestra il training distribuito. Ogni agente \`e mappato a una policy PPO indipendente tramite il \texttt{policy\_mapping\_fn}.
    \item \textbf{Dati PVGIS}: profili di irradianza solare oraria reali per Roma (lat.\ 41.893, lon.\ 12.483), anno 2023, con risoluzione oraria (8760 punti/anno).
\end{itemize}

\subsection{Topologia del Cluster: 8 Nodi, 4 Orientamenti}

Il cluster \`e composto da \textbf{8 nodi edge} organizzati in 4 coppie. Ogni coppia condivide lo stesso hardware \emph{e} lo stesso orientamento del pannello solare. Il mapping \`e progettato per \textbf{bilanciare} le capacit\`a hardware con la disponibilit\`a solare:

\begin{table}[h]
\centering
\caption{Mapping Nodi--Hardware--Orientamento}
\label{tab:mapping}
\begin{tabular}{@{}cccccc@{}}
\toprule
\textbf{Nodi} & \textbf{Hardware} & $P_{\text{compl}}$ & $B_{\max}$ (Wh) & \textbf{Orientamento} & \textbf{Logica} \\
\midrule
$W_0, W_1$ & Jetson Orin   & 0.95 & 10.0 & Nord  & Grande batteria compensa poco sole \\
$W_2, W_3$ & Jetson Nano   & 0.80 & 8.0  & Ovest & Sole pomeridiano \\
$W_4, W_5$ & Raspberry Pi  & 0.65 & 7.0  & Est   & Sole mattutino \\
$W_6, W_7$ & Low-power MCU & 0.50 & 5.0  & Sud   & Molto sole compensa batteria piccola \\
\bottomrule
\end{tabular}
\end{table}

La scelta di assegnare i nodi pi\`u deboli (batteria da 5\,Wh) all'esposizione pi\`u favorevole (Sud) e i nodi pi\`u potenti (batteria da 10\,Wh) all'esposizione peggiore (Nord) crea un sistema intrinsecamente bilanciato, dove nessuna coppia ha un vantaggio strutturale dominante.

% ======================================================================
\section{Modello dell'Ambiente}
% ======================================================================

\subsection{Spazio delle Osservazioni}

Ogni agente $i$ riceve ad ogni step $t$ un vettore di osservazione a 5 dimensioni:

\begin{equation}
\mathbf{o}_i^{(t)} = \Big[\; \underbrace{\tau^{(t)}}_{\text{tipo task}},\;\; \underbrace{q_i^{(t)}}_{\text{coda locale}},\;\; \underbrace{\hat{b}_i^{(t)}}_{\text{batteria norm.}},\;\; \underbrace{\bar{b}^{(t)}}_{\text{media cluster}},\;\; \underbrace{s_i^{(t)}}_{\text{sole attuale}} \;\Big]
\label{eq:obs}
\end{equation}

dove:
\begin{itemize}[nosep]
    \item $\tau^{(t)} \in \{0, 1, 2\}$: tipo del task in arrivo (60fps, 30fps, 15fps);
    \item $q_i^{(t)} \in \{0, \ldots, K\}$: lunghezza della coda locale ($K=5$);
    \item $\hat{b}_i^{(t)} = b_i^{(t)} / B_i^{\max} \in [0, 1]$: batteria normalizzata;
    \item $\bar{b}^{(t)} = \frac{1}{N}\sum_{j=1}^{N} \hat{b}_j^{(t)}$: media delle batterie normalizzate del cluster;
    \item $s_i^{(t)} \in [0, 1]$: irradianza solare normalizzata per l'orientamento del nodo $i$ all'istante $t$.
\end{itemize}

\textbf{Nota chiave}: l'agente vede sia la propria batteria $\hat{b}_i$ sia la media del cluster $\bar{b}$. Questo gli fornisce il segnale per capire se \`e \emph{sopra} o \emph{sotto} la media, senza necessitare di comunicazione esplicita.

\subsection{Spazio delle Azioni}

Lo spazio delle azioni \`e \textbf{discreto binario}:
\begin{equation}
a_i^{(t)} \in \{0, 1\} \quad \text{dove} \quad
\begin{cases}
0 & \text{Rifiuta il task} \\
1 & \text{Accetta il task (volontariato)}
\end{cases}
\end{equation}

\subsection{Logica del Dispatcher}

Quando pi\`u agenti scelgono $a_i = 1$ (volontari), il \textbf{dispatcher} assegna il task secondo questa logica:
\begin{enumerate}[nosep]
    \item Seleziona i volontari con la coda pi\`u corta: $\mathcal{V}^* = \arg\min_{j \in \mathcal{V}} q_j$;
    \item In caso di parit\`a, sceglie casualmente tra i candidati;
    \item Se la coda del selezionato \`e piena ($q_j = K$), il task viene \textbf{rifiutato};
    \item Se nessuno si offre volontario e il cloud \`e disponibile, il task va al \textbf{cloud}; altrimenti viene \textbf{rifiutato}.
\end{enumerate}

\subsection{Dinamica della Batteria}

Ad ogni step, la batteria di ogni nodo $i$ si aggiorna secondo:

\begin{equation}
b_i^{(t+1)} = \min\!\Big(\; B_i^{\max},\;\; b_i^{(t)} - c_{\text{idle}} - c_{\text{work}} \cdot \mathbb{1}[\text{task completato}] + s_i^{(t)} \cdot R_{\max} \;\Big)
\label{eq:battery}
\end{equation}

dove:
\begin{itemize}[nosep]
    \item $c_{\text{idle}} = 0.001$\,Wh: consumo base per step (sempre presente);
    \item $c_{\text{work}} = 0.05$\,Wh: consumo aggiuntivo per completamento task;
    \item $R_{\max} = 0.06$\,Wh: ricarica massima per step a pieno sole;
    \item $\mathbb{1}[\text{task completato}]$: indicatore stocastico, dipende da $P_{\text{compl},i}$.
\end{itemize}

Il completamento del task \`e stocastico: ad ogni step, se la coda \`e non vuota, il task viene completato con probabilit\`a $P_{\text{compl},i}$ (la velocit\`a del nodo dalla Tabella~\ref{tab:mapping}).

\textbf{Bilancio energetico}: a pieno sole ($s_i = 1$), un nodo che lavora consuma $0.05 + 0.001 = 0.051$\,Wh e ricarica $0.06$\,Wh, con un guadagno netto di $+0.009$\,Wh/step. Un nodo idle a pieno sole guadagna $0.06 - 0.001 = +0.059$\,Wh/step. Di notte ($s_i = 0$), il nodo perde $0.001$\,Wh/step in idle e $0.051$\,Wh/step se lavora.

% ======================================================================
\section{Dati Solari}
% ======================================================================

\subsection{Fonte e Formato}

I profili solari provengono dal database \textbf{PVGIS-SARAH3} (Photovoltaic Geographical Information System) della Commissione Europea. Parametri della query:
\begin{itemize}[nosep]
    \item Localit\`a: Roma (41.893\textdegree N, 12.483\textdegree E), elevazione 42\,m;
    \item Inclinazione pannello: 30\textdegree;
    \item Tecnologia: Silicio cristallino (c-Si), 1\,kWp nominale;
    \item Perdite di sistema: 14\%;
    \item 4 azimut: Sud (0\textdegree), Est ($-90$\textdegree), Ovest ($+90$\textdegree), Nord ($-179$\textdegree).
\end{itemize}

Ogni file CSV contiene 8760 valori orari della potenza $P$ (W) prodotta dal pannello per l'intero anno 2023.

\subsection{Normalizzazione}

La potenza \`e normalizzata dinamicamente rispetto al picco massimo presente nei dati:
\begin{equation}
s_{\text{orient}}^{(t)} = \text{clip}\!\left(\frac{P_{\text{orient}}^{(t)}}{P_{\max}},\; 0,\; 1\right), \quad P_{\max} = \max_{\text{orient}, t}\; P_{\text{orient}}^{(t)}
\end{equation}

Questo garantisce che il range $[0, 1]$ dell'observation space sia pienamente utilizzato.

\subsection{Data Augmentation (Solo Training)}

Durante il training, ogni episodio applica un fattore di augmentation $\xi \sim \mathcal{U}(0.5, 1.0)$ che moltiplica tutti i valori solari:
\begin{equation}
\tilde{s}_i^{(t)} = s_i^{(t)} \cdot \xi
\end{equation}

Questo simula condizioni meteo variabili (nuvolosit\`a parziale), rendendo la policy robusta a giornate con irradianza ridotta rispetto ai dati storici. In fase di valutazione, $\xi = 1$ (dati reali non alterati).

% ======================================================================
\section{Funzione di Reward}
% ======================================================================

La reward \`e \textbf{per-agente} (ogni nodo riceve una reward potenzialmente diversa) e si compone di due termini bilanciati dal parametro $\alpha = 0.5$:

\begin{equation}
r_i^{(t)} = \alpha \cdot r_{\text{fps}}^{(t)} + (1 - \alpha) \cdot r_{\text{batt},i}^{(t)}
\label{eq:reward}
\end{equation}

\subsection{Componente Prestazionale $r_{\text{fps}}$ (Condivisa)}

Questa componente \`e identica per tutti gli agenti e premia/penalizza l'esito del task:

\begin{equation}
r_{\text{fps}} =
\begin{cases}
-2.0 & \text{task rifiutato (perso per il sistema)} \\
-1.5 & \text{task al cloud, ma c'era spazio edge} \\
+0.5 & \text{task al cloud, edge pieno (salvavita)} \\
+0.1 & \text{task edge tipo 0 (leggero)} \\
+1.0 & \text{task edge tipo 1 (medio)} \\
+2.0 & \text{task edge tipo 2 (pesante)}
\end{cases}
\end{equation}

\subsection{Componente Energetica $r_{\text{batt},i}$ (Individuale)}

Questa componente contiene tre sotto-termini che risolvono il problema del \emph{credit assignment}:

\paragraph{(a) Penalit\`a sulla Varianza (condivisa).} Quando un task viene accettato sull'edge:
\begin{equation}
r_{\text{var}} = -20 \cdot \text{Var}\!\left(\hat{b}_1, \ldots, \hat{b}_N\right) = -20 \cdot \frac{1}{N}\sum_{j=1}^{N}\left(\hat{b}_j - \bar{b}\right)^2
\end{equation}

Questo termine \`e il cuore del bilanciamento: qualsiasi azione che aumenta la dispersione delle batterie normalizzate viene penalizzata.

\paragraph{(b) Deviazione Individuale.} Penalizza specificamente l'agente $i$ per quanto si discosta dalla media:
\begin{equation}
r_{\text{dev},i} = -5 \cdot \left|\hat{b}_i - \bar{b}\right|
\end{equation}

Questo termine risolve il credit assignment: un agente con batteria molto diversa dalla media riceve una penalità maggiore, indipendentemente dalle azioni degli altri.

\paragraph{(c) Anti-Laziness.} Penalizza gli agenti che rifiutano il task ($a_i = 0$) pur avendo batteria sopra la media \emph{e} sole disponibile:
\begin{equation}
r_{\text{lazy},i} =
\begin{cases}
-0.5 \cdot s_i^{(t)} & \text{se } a_i = 0,\;\; \hat{b}_i > \bar{b},\;\; s_i^{(t)} > 0.2,\;\; \text{task non rifiutato} \\
0 & \text{altrimenti}
\end{cases}
\end{equation}

La combinazione per un task accettato sull'edge \`e quindi:
\begin{equation}
r_{\text{batt},i} = r_{\text{var}} + r_{\text{dev},i} + r_{\text{lazy},i}
\end{equation}

Quando il task \`e rifiutato o inviato al cloud, $r_{\text{batt},i}$ usa penalità fisse basate sul tipo di task (da 0.0 a $-0.4$), pi\`u l'eventuale termine anti-laziness.

% ======================================================================
\section{Training IPPO}
% ======================================================================

\subsection{Paradigma: Independent PPO}

Ogni nodo $W_i$ possiede una \textbf{policy neurale indipendente} $\pi_{\theta_i}$. Non c'\`e condivisione di pesi, comunicazione tra agenti, o critic centralizzato. Ogni policy \`e una rete neurale standard di PPO che mappa osservazioni locali $\mathbf{o}_i$ in una distribuzione sull'azione $a_i$.

Le 8 policy vengono addestrate \textbf{simultaneamente} nello stesso ambiente: ad ogni step, ciascuna osserva il proprio stato, decide, e riceve la propria reward. L'ambiente funge da unico canale di coordinamento implicito (attraverso $\bar{b}$ e la varianza).

\subsection{Iperparametri}

\begin{table}[h]
\centering
\caption{Configurazione PPO}
\label{tab:hyperparams}
\begin{tabular}{@{}ll@{}}
\toprule
\textbf{Parametro} & \textbf{Valore} \\
\midrule
Algoritmo & PPO (via Ray RLlib) \\
Discount factor $\gamma$ & 0.99 \\
Learning rate & $5 \times 10^{-4}$ \\
Train batch size & 8000 \\
Minibatch size & 256 \\
Num.\ env runners & 1 \\
Iterazioni totali & 1000 \\
Durata episodio & 10000 step (max) \\
\bottomrule
\end{tabular}
\end{table}

\subsection{Strategie di Robustezza}

Per evitare che la policy memorizzi pattern specifici del dataset:

\begin{enumerate}[nosep]
    \item \textbf{Punto di partenza casuale}: ad ogni reset, il \texttt{weather\_step} iniziale \`e campionato uniformemente nell'anno, evitando che tutti gli episodi inizino da gennaio.
    \item \textbf{Batterie casuali}: in training, le batterie iniziali sono campionate da $\mathcal{U}(0.4, 1.0) \cdot B_i^{\max}$, forzando la policy a gestire scenari di batteria bassa.
    \item \textbf{Augmentation solare}: il fattore $\xi \sim \mathcal{U}(0.5, 1.0)$ riduce stocasticamente l'irradianza, simulando giornate nuvolose non presenti nei dati storici.
\end{enumerate}

% ======================================================================
\section{Comportamento Atteso: la Solar Wave}
% ======================================================================

Il comportamento emergente desiderato \`e il \textbf{Load Shifting guidato dal sole}:

\begin{enumerate}[nosep]
    \item \textbf{Mattina} (ore 6--12): i nodi Est ($W_4, W_5$) ricevono il picco solare. Dovrebbero accettare pi\`u task, sfruttando la ricarica per compensare il consumo.
    \item \textbf{Mezzogiorno} (ore 11--15): i nodi Sud ($W_6, W_7$) raggiungono il picco. Nonostante l'hardware debole, il sole abbondante consente loro di lavorare attivamente.
    \item \textbf{Pomeriggio} (ore 13--18): i nodi Ovest ($W_2, W_3$) entrano nel loro picco. Il carico si sposta su di essi.
    \item \textbf{Notte} (ore 18--6): tutti i nodi riducono l'accettazione al minimo. I nodi Nord ($W_0, W_1$), con batterie grandi, possono gestire i pochi task notturni essenziali.
\end{enumerate}

Questo comportamento non \`e programmato esplicitamente, ma dovrebbe \textbf{emergere} dall'interazione tra:
\begin{itemize}[nosep]
    \item L'osservazione del sole locale $s_i^{(t)}$ (l'agente sa quanto sole ha);
    \item La penalità sulla varianza (accettare senza sole scarica la batteria e aumenta la varianza);
    \item Il termine anti-laziness (rifiutare con sole e batteria alta viene penalizzato);
    \item La deviazione individuale (chi si allontana dalla media viene penalizzato di pi\`u).
\end{itemize}

% ======================================================================
\section{Metriche di Valutazione}
% ======================================================================

La valutazione del sistema addestrato avviene in modalit\`a deterministica (\texttt{explore=False}), con batterie inizializzate al 100\% e dati solari reali ($\xi = 1$). Le metriche chiave sono:

\subsection{Varianza Energetica del Cluster}

\begin{equation}
V^{(t)} = \text{Var}\!\left(\hat{b}_1^{(t)}, \ldots, \hat{b}_N^{(t)}\right)
\end{equation}

Tracciata nel tempo, indica quanto il sistema riesce a mantenere bilanciate le batterie. Il target ideale \`e $V^{(t)} \approx 0$ per tutta la simulazione.

\subsection{Task Rifiutati}

Il conteggio totale dei task rifiutati durante la simulazione. Un sistema ben addestrato dovrebbe minimizzare i rifiuti, accettando i task sui nodi con pi\`u energia disponibile.

\subsection{Profili di Batteria}

I livelli di batteria normalizzati di tutti gli 8 nodi, tracciati nel tempo. Il comportamento atteso \`e un andamento ``a onda'': i nodi con sole si caricano mentre lavorano, i nodi senza sole si scaricano lentamente in idle, con le curve che si avvicinano e si allontanano seguendo il ciclo solare.

\subsection{Lunghezza Media delle Code}

La media delle code di tutti i nodi, tracciata nel tempo. Indica l'efficienza di smaltimento dei task e la distribuzione del carico.

\subsection{Protocollo di Test}

\begin{enumerate}[nosep]
    \item Simulazione di \textbf{168 step} (1 settimana) a partire da un seed fisso (\texttt{seed=42}) per riproducibilit\`a;
    \item Generazione di 3 grafici: (1) Batterie degli 8 nodi, (2) Varianza di sistema, (3) Code medie + task rifiutati;
    \item Confronto qualitativo con una \textbf{Random Policy} (baseline) per dimostrare che il training ha prodotto un comportamento intelligente.
\end{enumerate}

% ======================================================================
\section{Riepilogo dei File del Progetto}
% ======================================================================

\begin{table}[h]
\centering
\caption{Struttura del progetto}
\begin{tabularx}{\textwidth}{@{}lX@{}}
\toprule
\textbf{File} & \textbf{Descrizione} \\
\midrule
\texttt{config.py} & Parametri globali: hardware, batterie, velocit\`a, iperparametri RL \\
\texttt{env/solar\_edge\_env.py} & Ambiente PettingZoo con dinamica batterie, ricarica solare, reward per-agente \\
\texttt{train\_solar\_ippo.py} & Script di training: configura RLlib, definisce 8 policy indipendenti, lancia PPO \\
\texttt{evaluate\_solar\_model.py} & Valutazione post-training: simulazione deterministica + generazione grafici \\
\texttt{solar\_data/*.csv} & 4 profili PVGIS (Sud, Est, Ovest, Nord) con 8760 valori orari ciascuno \\
\bottomrule
\end{tabularx}
\end{table}

\end{document}
